\chapter{Lineaire differentiaalvergelijkingen}\label{ch:b1:diffeq}

Differentiaalvergelijkingen kwamen voor het eerst voor in het
bovenbouwvolume; hier worden ze met volledige bewijzen en in grotere
algemeenheid behandeld: lineaire vergelijkingen van de eerste orde met
variabele coëfficiënten (volledig opgelost met de methode van de \omterm{thm:b1:diffeq:voc}{variatie
van constanten}) en lineaire vergelijkingen van de tweede orde met
constante coëfficiënten, het model voor oscillaties. Beide gevallen
vertonen dezelfde structuur: \emph{algemene oplossing $=$ één particuliere
oplossing $+$ algemene oplossing van de homogene vergelijking}.

\section{Lineaire vergelijkingen van de eerste orde}

\begin{definition}\label{def:b1:diffeq:linear1}
Zij $I$ een interval en zij $a, b \colon I \to \R$ (of $\C$) continu. De
vergelijking
\[
(E)\colon\quad y' + a(x)\,y = b(x),
\]
in de onbekende afleidbare functie $y \colon I \to \R$ (of $\C$), heet een
\emph{lineaire differentiaalvergelijking van de eerste
orde}\index{differentiaalvergelijking!lineair van de eerste orde}. De
vergelijking $(H)\colon y' + a(x) y = 0$ is haar \emph{homogene}
vergelijking.
\end{definition}

\begin{theorem}[De homogene vergelijking oplossen]\label{thm:b1:diffeq:homogeneous1}
Zij $A$ een primitieve functie van $a$ op $I$ (die bestaat:
\cref{ch:b1:integration}). De oplossingen van $(H)$ op $I$ zijn precies de
functies
\[
y(x) = \lambda\, \eu^{-A(x)}, \qquad \lambda \in \R \text{ (of } \C).
\]
\end{theorem}

\begin{proof}
Deze functies zijn oplossingen: $y' = -\lambda A' \eu^{-A} = -a y$.
Omgekeerd, zij $y$ een oplossing van $(H)$ en zet $z(x) = y(x)\,
\eu^{A(x)}$. Dan is
\[
z' = y' \eu^{A} + y\, a\, \eu^{A} = (y' + ay)\, \eu^{A} = 0 ,
\]
zodat $z$ constant is op het interval $I$, zeg $z = \lambda$: $y = \lambda
\eu^{-A}$. (Let op de logica: er gaat geen oplossing verloren, want
\emph{elke} oplossing is in de aangekondigde vorm geschreven.)
\end{proof}

\begin{example}[Een homogene vergelijking met variabele coëfficiënt]\label{ex:b1:diffeq:varcoeff}
Los $y' + (\cos x)\,y = 0$ op op $\R$. Een primitieve functie van $a(x)
= \cos x$ is $A(x) = \sin x$, dus de oplossingen zijn
\[
y(x) = \lambda\,\eu^{-\sin x}, \qquad \lambda \in \R .
\]
Twee lezingen. Elke oplossing is periodiek (periode $2\pi$) en wordt
nooit nul tenzij $\lambda = 0$ --- het teken van $\lambda$ is voor
altijd het teken van $y$, want een exponentiële kan de nul niet
oversteken. En de oplossing door $y(0) = y_0$ is $y_0\eu^{-\sin x}$:
precies één kromme van de familie door elk beginpunt, het
eendimensionale beeld van \cref{thm:b1:diffeq:voc}~(2).
\end{example}

\begin{theorem}[Variatie van constanten; Cauchy-probleem]\label{thm:b1:diffeq:voc}
Met bovenstaande notaties:
\begin{enumerate}
  \item De oplossingen van $(E)$ op $I$ zijn precies
        \[
        y(x) = \Bigl(\lambda + \int_{x_0}^{x} b(t)\, \eu^{A(t)} \dd t
        \Bigr)\, \eu^{-A(x)},
        \qquad \lambda \in \R,
        \]
        met $x_0 \in I$ vast. Equivalent: algemene oplossing van $(H)$
        plus één particuliere oplossing van $(E)$.
  \item Voor elke $x_0 \in I$ en elke $y_0$ heeft het
        \emph{Cauchy-probleem}\index{Cauchy-probleem} ``$(E)$ en
        $y(x_0) = y_0$'' precies één oplossing op $I$.
\end{enumerate}
\end{theorem}

\begin{proof}
(1) Volgens de methode die \emph{variatie van
constanten}\index{variatie van constanten} heet, zoeken we oplossingen
van de vorm $y = \mu(x)\, \eu^{-A(x)}$ met $\mu$ afleidbaar --- daarbij
gaat geen algemeenheid verloren, want elke functie op $I$ laat zich zo
schrijven ($\mu = y\,\eu^{A}$). Substitutie geeft
\[
y' + ay = \mu' \eu^{-A} - \mu a \eu^{-A} + a\mu\eu^{-A}
= \mu'\, \eu^{-A},
\]
zodat $y$ een oplossing van $(E)$ is dan en slechts dan als $\mu'(x) =
b(x)\,\eu^{A(x)}$, dan en slechts dan als $\mu(x) = \lambda +
\int_{x_0}^x b(t)\eu^{A(t)}\dd t$ voor een zekere constante $\lambda$
(twee primitieve functies van dezelfde continue functie op een interval
verschillen een constante).

(2) In de formule is $y(x_0) = \lambda\,\eu^{-A(x_0)}$: de voorwaarde
$y(x_0) = y_0$ legt $\lambda = y_0 \eu^{A(x_0)}$ ondubbelzinnig vast.
\end{proof}

\begin{example}\label{ex:b1:diffeq:order1}
Los $y' + \dfrac{y}{x} = x^2$ op op $I = \intoo{0}{+\infty}$. Hier is
$a(x) = \frac 1x$, $A(x) = \ln x$ en $\eu^{-A(x)} = \frac 1x$. Homogene
oplossingen: $\frac{\lambda}{x}$. \omterm{thm:b1:diffeq:voc}{Variatie van constanten}: $\mu'(x) = x^2
\cdot x = x^3$, dus $\mu = \frac{x^4}{4} + \lambda$, en
\[
y(x) = \frac{x^3}{4} + \frac{\lambda}{x}, \qquad \lambda \in \R .
\]
Met de beginvoorwaarde $y(1) = 0$ is $\lambda = -\frac14$.
\emph{Controle:} $y' + \frac yx = \frac{3x^2}{4} - \frac{\lambda}{x^2} +
\frac{x^2}{4} + \frac{\lambda}{x^2} = x^2$.
\end{example}

\begin{example}[Raden verslaat integreren]\label{ex:b1:diffeq:guess}
Los $y' + 2x\,y = x$ op op $\R$. \omterm{thm:b1:diffeq:voc}{Variatie van constanten} werkt ($A =
x^2$, $\mu' = x\,\eu^{x^2}$, $\mu = \frac12\eu^{x^2} + \lambda$), maar
opmerken dat de \emph{constante} $y_p = \frac12$ de vergelijking oplost
($0 + 2x\cdot\frac12 = x$) gaat sneller. Met de homogene oplossingen
$\lambda\,\eu^{-x^2}$:
\[
y(x) = \frac12 + \lambda\,\eu^{-x^2}, \qquad \lambda \in \R .
\]
Elke oplossing convergeert bijzonder snel naar $\frac12$ als $x \to
\pm\infty$: de constante particuliere oplossing is een
\emph{evenwicht} waar alle oplossingen bij aansluiten. Het inzicht: vóór
je de algemene methode start, besteed je tien seconden aan het zoeken
naar een voor de hand liggende particuliere oplossing (een constante, een
monoom, een veelvoud van het rechterlid); de structuurstelling maakt het
werk dan af.
\end{example}

\begin{remark}[Intervallen doen ertoe]
De stelling leeft op een \emph{interval} waar $a$ en $b$ continu zijn.
Voor $y' + \frac yx = 0$ op $\R^*$ zijn de oplossingen
$\frac{\lambda}{x}$ op $\intoo{0}{+\infty}$ en $\frac{\mu}{x}$ op
$\intoo{-\infty}{0}$ met \emph{onafhankelijke} constanten: er is geen
reden waarom één formule over de singulariteit in $0$ heen zou lijmen.
\end{remark}

\begin{example}[Eén complex rechterlid, twee reële antwoorden]\label{ex:b1:diffeq:complexrhs}
Los $y' - y = \cos x$ en $y' - y = \sin x$ in één beweging op. Werk in
$\C$ met het rechterlid $\eu^{\iu x}$: de poging $y_p = c\,\eu^{\iu x}$
geeft $c(\iu - 1)\eu^{\iu x} = \eu^{\iu x}$, dus
\[
c = \frac1{\iu - 1} = \frac{-1 - \iu}2,
\qquad
y_p = -\frac{(1 + \iu)(\cos x + \iu\sin x)}2
= \frac{\sin x - \cos x}2 + \iu\,\frac{-\sin x - \cos x}2 .
\]
Omdat de vergelijking reële coëfficiënten heeft, splitsen reëel en
imaginair deel: $\frac{\sin x - \cos x}2$ lost $y' - y = \cos x$ op, en
$-\frac{\sin x + \cos x}2$ lost $y' - y = \sin x$ op (controleer de
eerste: de afgeleide $\frac{\cos x + \sin x}2$ min de functie geeft
$\cos x$). Eén complexe regel verving twee keer \omterm{thm:b1:diffeq:voc}{variatie van constanten}
--- dezelfde zuinigheid die \cref{met:b1:diffeq:particular} voor de
tweede orde systematiseert, en een terugkerend dividend van
\cref{ch:b1:complex}.
\end{example}

\section{Lineaire vergelijkingen van de tweede orde met constante coëfficiënten}

\begin{definition}\label{def:b1:diffeq:linear2}
Zij $a, b \in \R$ en zij $f \colon I \to \R$ continu. De vergelijking
\[
(E)\colon\quad y'' + a\,y' + b\,y = f(x)
\]
heet een \emph{lineaire vergelijking van de tweede orde met constante
coëfficiënten}\index{differentiaalvergelijking!lineair van de tweede
orde}; $(H)\colon y'' + ay' + by = 0$ is haar homogene vergelijking en
$\chi(r) = r^2 + ar + b$ haar \emph{karakteristieke
veelterm}\index{karakteristieke veelterm}.
\end{definition}

\begin{theorem}[Homogene oplossingen]\label{thm:b1:diffeq:homogeneous2}
Zij $\Delta = a^2 - 4b$ de discriminant van $\chi$. De reële oplossingen
van $(H)$ op $\R$ zijn:
\begin{enumerate}
  \item als $\Delta > 0$, met $r_1 \neq r_2$ de twee reële wortels:
        $\;y = \lambda\, \eu^{r_1 x} + \mu\, \eu^{r_2 x}$;
  \item als $\Delta = 0$, met $r_0$ de dubbele wortel:
        $\;y = (\lambda + \mu x)\, \eu^{r_0 x}$;
  \item als $\Delta < 0$, met wortels $\alpha \pm \iu\omega$
        ($\omega > 0$):
        $\;y = \eu^{\alpha x} (\lambda \cos\omega x + \mu
        \sin\omega x)$;
\end{enumerate}
telkens met $(\lambda, \mu)$ dat $\R^2$ doorloopt.
\end{theorem}

\begin{proof}
Merk eerst op dat voor $r \in \C$ de functie $x \mapsto \eu^{rx}$ een
oplossing van $(H)$ is dan en slechts dan als $\chi(r) = 0$
(substitutie: $(r^2 + ar + b)\eu^{rx} = 0$). Daarom zijn exponentiëlen de
natuurlijke eerste gok: differentiëren werkt op $\eu^{rx}$ als
vermenigvuldigen met het getal $r$, zodat de differentiaalvergelijking
de numerieke vergelijking $\chi(r) = 0$ wordt --- het hele analytische
probleem wordt samengeperst tot het vinden van de wortels van één
vierkantsvergelijking.

De sleutelstap is een verandering van onbekende die de orde verlaagt. Zij
$r$ een (eventueel complexe) wortel van $\chi$ en schrijf $y = z\,
\eu^{rx}$, waarbij geen algemeenheid verloren gaat. Dan is
\[
y'' + ay' + by
= \bigl(z'' + (2r + a) z' + \chi(r) z\bigr)\eu^{rx}
= \bigl(z'' + (2r + a) z'\bigr)\eu^{rx},
\]
zodat $(H)$ de vergelijking $u' + (2r + a) u = 0$ van de \emph{eerste
orde} wordt voor $u = z'$.

\emph{Geval $\Delta \neq 0$:} kies $r = r_1$; dan is $2r_1 + a = r_1 -
r_2$ (want $r_1 + r_2 = -a$). Volgens
\cref{thm:b1:diffeq:homogeneous1} is $z' = c\,\eu^{(r_2 - r_1)x}$ voor
een zekere constante $c$; integreren op $\R$ geeft $z = \mu\, \eu^{(r_2 -
r_1)x} + \lambda$ met $\mu = \frac{c}{r_2 - r_1}$, en dus $y = z\,
\eu^{r_1 x} = \lambda\,\eu^{r_1x} + \mu\,\eu^{r_2x}$. Is $\Delta < 0$,
dan zijn de wortels $\alpha \pm \iu\omega$ en zijn de complexe
oplossingen $y = c_1\eu^{(\alpha+\iu\omega)x} +
c_2\eu^{(\alpha-\iu\omega)x}$ met $c_1, c_2 \in \C$. Welke daarvan zijn
reëelwaardig? Omdat $\conj{\eu^{(\alpha+\iu\omega)x}} =
\eu^{(\alpha-\iu\omega)x}$, is de toegevoegde van $y$ gelijk aan
$\conj{c_2}\, \eu^{(\alpha+\iu\omega)x} +
\conj{c_1}\,\eu^{(\alpha-\iu\omega)x}$, en $y = \conj y$ voor alle $x$
dwingt $c_2 = \conj{c_1}$ af (de twee exponentiëlen zijn lineair
onafhankelijk: evalueer in twee punten, of vergelijk in $x=0$ na deling
door $\eu^{\alpha x}$). Schrijven we $c_1 = \frac{\lambda - \iu\mu}2$ met
reële $\lambda, \mu$, dan is
\[
y = 2\,\Re\Bigl(c_1\,\eu^{\alpha x}(\cos\omega x +
\iu\sin\omega x)\Bigr)
= \eu^{\alpha x}\bigl(\lambda\cos\omega x + \mu\sin\omega
x\bigr),
\]
en omgekeerd is elke zo'n functie een oplossing (reëel deel van een
complexe oplossing van een reële vergelijking): de reële
oplossingsruimte is zoals aangekondigd.

\emph{Geval $\Delta = 0$:} $r = r_0$ en $2r_0 + a = 0$, dus $z'' = 0$: $z
= \lambda + \mu x$ en $y = (\lambda + \mu x)\eu^{r_0 x}$.
\end{proof}

\begin{example}[Een Cauchy-probleem van begin tot eind]\label{ex:b1:diffeq:cauchy2}
Los $y'' - 3y' + 2y = 0$ op met $y(0) = 0$ en $y'(0) = 1$. De
\omterm{def:b1:diffeq:linear2}{karakteristieke veelterm} $r^2 - 3r + 2 = (r - 1)(r - 2)$ heeft de
reële wortels $1$ en $2$: de algemene oplossing is $y = \lambda\eu^{x} +
\mu\eu^{2x}$. De twee voorwaarden geven het lineaire stelsel
\[
\lambda + \mu = 0, \qquad \lambda + 2\mu = 1 ,
\]
dus $\mu = 1$ en $\lambda = -1$:
\[
y(x) = \eu^{2x} - \eu^{x} .
\]
Controle: $y(0) = 0$; $y' = 2\eu^{2x} - \eu^x$ geeft $y'(0) = 1$; en $y''
- 3y' + 2y = (4 - 6 + 2)\eu^{2x} + (-1 + 3 - 2)\eu^x = 0$. Let op de
vorm van het antwoord: bij $-\infty$ overheerst de trage mode $-\eu^x$,
bij $+\infty$ de snelle mode $\eu^{2x}$. Oplossingen lezen als
superposities van modes met verschillende afname- of groeisnelheden is de
lonende gewoonte --- zo is ook de splitsing tussen inschakelverschijnsel
en regime in de weekendopgave georganiseerd.
\end{example}

\begin{omfigure}
\begin{tikzpicture}[xscale=1.35, yscale=1.55]
  \draw[->, thin] (0,0) -- (6.4,0) node[below] {$x$};
  \draw[->, thin] (0,-0.85) -- (0,1.25) node[left] {$y$};
  \draw[omDef, thick, domain=0:6.2, samples=400, smooth]
    plot (\x, {exp(-0.35*\x)*cos(4*\x r)});
  \draw[omThm, thick, domain=0:6.2, samples=120, smooth]
    plot (\x, {(1+1.8*\x)*exp(-1.1*\x)});
  \draw[omProp, thick, domain=0:6.2, samples=120, smooth]
    plot (\x, {1.15*exp(-0.45*\x) - 0.15*exp(-2.5*\x)});
  \node[omDef, font=\small] at (2.1,-0.62)
    {$\Delta < 0$: oscilleert};
  \node[omThm, font=\small] at (0.95,1.1) {$\Delta = 0$};
  \node[omProp, font=\small] at (4.9,0.38)
    {$\Delta > 0$: geen oscillatie};
\end{tikzpicture}

{\small De drie regimes van $y'' + ay' + by = 0$ met afnemende
oplossingen: gedempte oscillatie (complexe wortels), kritieke terugkeer
(dubbele wortel), overgedempte afname (twee reële wortels). Welk regime
optreedt lees je af aan het teken van $\Delta = a^2 - 4b$ alleen ---
nog vóór je iets oplost.}
\end{omfigure}

\begin{theorem}[Structuur en Cauchy-probleem]\label{thm:b1:diffeq:structure2}
\begin{enumerate}
  \item Is $y_p$ één particuliere oplossing van $(E)$, dan zijn de
        oplossingen van $(E)$ precies $y_p + y_h$, waarbij $y_h$ de
        oplossingen van $(H)$ doorloopt.
  \item (Superpositie) Lost $y_1$ de vergelijking $y'' + ay' + by =
        f_1$ op en $y_2$ die met rechterlid $f_2$, dan lost $y_1 + y_2$
        de vergelijking met rechterlid $f_1 + f_2$ op.
  \item Voor alle $x_0 \in I$ en $(y_0, y_0')$ heeft het
        \omterm{thm:b1:diffeq:voc}{Cauchy-probleem} ``$(E)$, $y(x_0) = y_0$, $y'(x_0) = y_0'$''
        precies één oplossing op $I$. \emph{(Bestaan onder voorbehoud
        van een particuliere oplossing; uniciteit volledig.)}
\end{enumerate}
\end{theorem}

\begin{proof}
(1) $y$ lost $(E)$ op precies wanneer $y - y_p$ de vergelijking $(H)$
oplost, wegens de lineariteit van $y \mapsto y'' + ay' + by$. (2) is
dezelfde lineariteit.

(3) Wegens (1) volstaat het te bewijzen dat de constanten $(\lambda,
\mu)$ altijd, en op één manier, aan willekeurige gegevens $(y_0, y_0')$
aangepast kunnen worden. Na verschuiving van de variabele mogen we $x_0 =
0$ aannemen. In geval (1) van
\cref{thm:b1:diffeq:homogeneous2} geeft $y = \lambda\eu^{r_1x} +
\mu\eu^{r_2x}$
\[
y(0) = \lambda + \mu, \qquad y'(0) = r_1\lambda + r_2\mu :
\]
een lineair stelsel in $(\lambda, \mu)$ met determinant $r_2 - r_1 \neq
0$; expliciet oplossen geeft $\mu = \frac{y_0' - r_1y_0}{r_2 - r_1}$ en
$\lambda = y_0 - \mu$: precies één oplossing. In geval (2) is $y(0) =
\lambda$ en $y'(0) = r_0\lambda + \mu$: het stelsel is driehoekig met
determinant $1$, met oplossing $\lambda = y_0$ en $\mu = y_0' - r_0y_0$.
In geval (3) is $y(0) = \lambda$ en $y'(0) = \alpha\lambda + \omega\mu$:
determinant $\omega \neq 0$, met oplossing $\lambda = y_0$ en $\mu =
\frac{y_0' - \alpha y_0}\omega$. In elk geval is de \omterm{def:b1:logic:map}{afbeelding} $(\lambda,
\mu) \mapsto (y(x_0), y'(x_0))$ een lineaire bijectie --- de taal van
\cref{ch:b1:linmaps} zal dit gevalsonderzoek tot één zin samenpersen.
\end{proof}

\begin{method}[Particuliere oplossing bij $f(x) = P(x)\,\eu^{\gamma x}$]\label{met:b1:diffeq:particular}
Is het rechterlid $P(x)\,\eu^{\gamma x}$ met $P$ een veelterm en $\gamma
\in \R$ (dat dekt veeltermen, exponentiëlen, en via complexe $\gamma$ of
superpositie ook $\cos$ en $\sin$), zoek dan een particuliere oplossing
van de vorm
\[
y_p(x) = x^{m}\, Q(x)\, \eu^{\gamma x},
\qquad
m = \text{multipliciteit van } \gamma \text{ als wortel van } \chi
\ (m = 0, 1 \text{ of } 2),
\]
met $Q$ een veelterm van dezelfde graad als $P$, waarvan je de
coëfficiënten vindt door substitutie en gelijkstelling. Voor $f =
K\cos\omega x$ (of $\sin$) los je op met rechterlid $K\eu^{\iu\omega x}$
en neem je het reële (respectievelijk imaginaire) deel.
\end{method}

\begin{example}[Superpositie in actie]\label{ex:b1:diffeq:superposition}
Los $y'' - y = \eu^{x} + 4$ op op $\R$. Homogeen: $\chi(r) = r^2 - 1$ met
wortels $\pm1$, dus $y_h = \lambda\eu^x + \mu\eu^{-x}$. Splits het
rechterlid en behandel elk stuk met het methodekader. \emph{Stuk
$\eu^x$:} hier is $\gamma = 1$ een enkelvoudige wortel van $\chi$, dus
probeer $y_1 = c\,x\,\eu^x$: dan is $y_1'' - y_1 = c(x + 2)\eu^x -
cx\eu^x = 2c\,\eu^x$, wat $c = \frac12$ geeft. \emph{Stuk $4$:} $\gamma =
0$ is geen wortel; de constante $y_2 = -4$ voldoet. Met de superpositie
(\cref{thm:b1:diffeq:structure2}~(2)):
\[
y = \frac{x\,\eu^x}2 - 4 + \lambda\,\eu^x + \mu\,\eu^{-x},
\qquad (\lambda, \mu) \in \R^2 .
\]
Merk op hoe de twee stukken \emph{verschillende} vormen eisten ($m = 1$
tegenover $m = 0$): de multipliciteitstoets pas je op elke exponent
afzonderlijk toe, en dat is de hele reden om het rechterlid te splitsen
vóór je gaat raden.
\end{example}

\begin{example}[De multipliciteitsregel aan het werk]\label{ex:b1:diffeq:multiplicity}
Los $y'' + y' = x$ op op $\R$. Het rechterlid is $P(x)\eu^{0 \cdot x}$
met $P(x) = x$, en $\gamma = 0$ is een \emph{enkelvoudige} wortel van
$\chi(r) = r^2 + r = r(r + 1)$: dus $m = 1$, en de juiste gok is $y_p =
x\,(\alpha x + \beta) = \alpha x^2 + \beta x$, één graad hoger dan $P$.
Substitutie geeft
\[
y_p'' + y_p' = 2\alpha + (2\alpha x + \beta)
= 2\alpha x + (2\alpha + \beta) ,
\]
en gelijkstellen met $x$ levert $\alpha = \frac12$ en $\beta = -1$: $y_p
= \frac{x^2}2 - x$. Algemene oplossing: $y = \frac{x^2}2 - x + \lambda +
\mu\,\eu^{-x}$. Hadden we $y_p = \alpha x + \beta$ gegokt (met
verwaarlozing van de multipliciteit), dan zou substitutie $y_p'' + y_p' =
\alpha$ geven, een constante --- geen enkele keuze van $\alpha, \beta$
kan $x$ evenaren, en dat falen is structureel: constanten lossen de
homogene vergelijking al op en zijn dus onzichtbaar voor het linkerlid.
De factor $x^m$ bestaat juist om uit de homogene oplossingsruimte te
klimmen.
\end{example}

\begin{remark}[De verzekeringspolis van dertig seconden]\label{rem:b1:diffeq:check}
Elke opgeloste vergelijking in dit hoofdstuk eindigt met een controle
door substitutie, en dat is niet ter versiering. Een berekening aan een
differentiaalvergelijking rijgt vele kleine stappen aaneen (een
primitieve functie, een productregel, twee constanten), en één
tekenfout plant zich onzichtbaar voort; de eindformule terug in de
vergelijking invullen vangt er vrijwel alle op, tegen de kostprijs van
één keer differentiëren. Kweek de reflex in drie lagen: controleer de
\emph{particuliere oplossing} apart (het homogene deel valt toch weg),
controleer de \emph{beginvoorwaarden} op de volledige oplossing, en
controleer, wanneer er een parameter in het spel is, een \emph{ontaarde
waarde} (levert de formule voor algemene $\Omega$ het bekende antwoord in
$\Omega = 0$?). De gewoonte kost een halve minuut en zet ``waarschijnlijk
juist'' om in ``geverifieerd''.
\end{remark}

\begin{remark}[Veelgemaakte fouten]\label{rem:b1:diffeq:pitfalls}
\begin{enumerate}
  \item \emph{Normaliseer eerst.} De formules gaan ervan uit dat de
        vergelijking $y' + a(x)y = b(x)$ luidt --- met coëfficiënt $1$
        bij $y'$. Voor $xy' - 2y = x^3$ deel je eerst door $x$ (op een
        interval dat $0$ mijdt) vóór je $a$ en $b$ afleest, zoals in
        \cref{exo:b1:diffeq:2}.
  \item \emph{Eén constante per dimensie, pas op het eind vastgelegd.}
        De algemene oplossing van de eerste orde draagt één constante,
        die van de tweede orde twee; beginvoorwaarden leg je op aan de
        \emph{volledige} oplossing $y_p + y_h$, nooit aan $y_h$ alleen
        --- ze opleggen vóór het optellen van $y_p$ is de meest
        voorkomende structurele fout.
  \item \emph{Let op de multipliciteit.} Een gok voor de particuliere
        oplossing die de homogene vergelijking oplost, is onzichtbaar
        voor het linkerlid; de factor $x^m$ van
        \cref{met:b1:diffeq:particular} is niet vrijblijvend
        (\cref{ex:b1:diffeq:multiplicity}).
  \item \emph{Intervallen horen bij het antwoord.} Oplossingen leven op
        intervallen waar de coëfficiënten continu zijn; over een
        singulariteit heen lijmen kan valse constanten scheppen
        (\cref{exo:b1:diffeq:12}) of de uniciteit vernietigen. ``Los op
        op $\R^*$'' betekent twee onafhankelijke problemen.
\end{enumerate}
\end{remark}

\begin{example}[Aandrijving buiten resonantie]\label{ex:b1:diffeq:offresonance}
Los $y'' + 4y = \sin x$ op op $\R$. Eigenfrequentie $2$, aandrijffrequentie
$1$: omdat $\iu$ \emph{geen} wortel van $\chi(r) = r^2 + 4$ is, is de
multipliciteit $m = 0$ en volstaat een gewone sinusoïde. De poging $y_p =
\alpha\sin x$ (een cosinus is niet nodig: de vergelijking heeft geen
$y'$-term, en $\sin$ brengt weer $\sin$ voort) geeft
\[
y_p'' + 4y_p = -\alpha\sin x + 4\alpha\sin x = 3\alpha\sin x ,
\]
dus $\alpha = \frac13$, en de algemene oplossing is
\[
y = \frac{\sin x}3 + \lambda\cos 2x + \mu\sin 2x .
\]
Elke oplossing blijft begrensd: een superpositie van twee oscillaties met
de frequenties $1$ (opgelegd) en $2$ (eigen). Vergelijk met het volgende
voorbeeld, waar aandrijving \emph{op} de eigenfrequentie de vorm van het
antwoord zelf verandert.
\end{example}

\begin{example}[Een aangedreven oscillatie]\label{ex:b1:diffeq:oscillation}
Los $y'' + y = \cos x$ op met $y(0) = 0$ en $y'(0) = 0$.

\emph{Homogeen:} $\chi(r) = r^2 + 1$ met wortels $\pm\iu$: $y_h =
\lambda\cos x + \mu \sin x$.

\emph{Particulier:} het rechterlid is $\Re(\eu^{\iu x})$ met $\gamma =
\iu$ een enkelvoudige wortel van $\chi$: probeer $z_p = c\, x\, \eu^{\iu
x}$ met $c \in \C$. Dan is $z_p'' + z_p = c\,(2\iu)\eu^{\iu x}$, wat
gelijk is aan $\eu^{\iu x}$ voor $c = \frac{1}{2\iu} = -\frac\iu2$. Dus
$z_p = -\frac{\iu}{2} x (\cos x + \iu \sin x)$ en $y_p = \Re(z_p) =
\frac{x \sin x}{2}$.

\emph{Algemene oplossing:} $y = \frac{x\sin x}{2} + \lambda\cos x +
\mu\sin x$. Voorwaarden: $y(0) = \lambda = 0$; en $y' = \frac{\sin x + x\cos
x}{2} + \mu\cos x$, dus $y'(0) = \mu = 0$. Antwoord: $y = \frac{x\sin
x}{2}$ --- een oscillatie waarvan de amplitude lineair groeit: het
verschijnsel \emph{resonantie}\index{resonantie}, veroorzaakt door het
systeem op zijn eigenfrequentie aan te drijven.
\end{example}

\begin{omfigure}
\begin{tikzpicture}
\begin{axis}[omaxis, width=11cm, height=5.2cm,
    xmin=0, xmax=25, ymin=-13, ymax=13,
    xtick={0,5,10,15,20,25}, ytick=\empty]
  \addplot[omThm, thick, domain=0:25, samples=400] {x*sin(deg(x))/2};
  \addplot[dashed, gray, domain=0:25] {x/2};
  \addplot[dashed, gray, domain=0:25] {-x/2};
\end{axis}
\end{tikzpicture}

{\small \omterm{ex:b1:diffeq:oscillation}{Resonantie}: de oplossing $y = \frac{x \sin x}{2}$ van $y'' + y =
\cos x$ oscilleert tussen de rechten $y = \pm\frac x2$ (streepjeslijn),
met een steeds groeiende amplitude.}
\end{omfigure}

\begin{remark}[Tussenspel: lineariteit is een meetkunde]\label{rem:b1:diffeq:interlude}
Kijk nog eens naar de vorm van elke oplossingsverzameling in dit
hoofdstuk: één bijzondere oplossing plus een ruimte homogene oplossingen
met één vrije constante (eerste orde) of twee (tweede orde). De
hoofdstukken over lineaire algebra
(\cref{ch:b1:vspaces,ch:b1:findim,ch:b1:linmaps}) zullen het exacte
vocabulaire leveren: de \omterm{def:b1:logic:map}{afbeelding} $L(y) = y'' + ay' + by$ is
\emph{lineair}, haar homogene oplossingen vormen de \emph{kern} van $L$,
een vectorruimte waarvan de \emph{dimensie} gelijk is aan de orde van de
vergelijking --- dat is de eerlijke inhoud van ``één constante per
orde'' --- en de oplossingsverzameling van $L(y) = f$ is een
\emph{affiene deelruimte}, een verschuiving van de kern. Zelfs de
\omterm{def:b1:logic:map}{Cauchy-afbeelding} $(\lambda, \mu) \mapsto (y(x_0), y'(x_0))$ van
\cref{thm:b1:diffeq:structure2} is een lineaire bijectie tussen twee
vlakken, dat wil zeggen een inverteerbaar $2 \times 2$-stelsel
(\cref{ch:b1:matrices}). Niets in dit hoofdstuk zal overgedaan hoeven
worden --- het zal alleen hernoemd worden, en dat hernoemen is de best
denkbare opwarming voor de lineaire algebra: elke abstracte definitie
daar heeft hier haar brood al verdiend.
\end{remark}

\begin{remark}[Waar dit hoofdstuk gebruikt wordt]\label{rem:b1:diffeq:whereused}
De structuurstelling --- de oplossingen van $(E)$ vormen ``één
particuliere oplossing plus de oplossingen van $(H)$'' --- is de eerste
verschijning van een patroon dat
\cref{ch:b1:vspaces,ch:b1:linmaps} zullen benoemen: de
oplossingsverzameling van $(H)$ is de \emph{kern} van de lineaire
\omterm{def:b1:logic:map}{afbeelding} $y \mapsto y'' + ay' + by$, en die van $(E)$ is een affiene
verschuiving ervan. De \omterm{def:b1:diffeq:linear2}{karakteristieke veelterm} keert terug als de
\omterm{def:b1:diffeq:linear2}{karakteristieke veelterm} van een matrix in \cref{ch:b1:matrices}: een
vergelijking van de tweede orde is een vermomd $2 \times 2$-stelsel van
de eerste orde, een standpunt dat het volume van bachelorjaar 2
systematiseert. De integralen die de \omterm{thm:b1:diffeq:voc}{variatie van constanten} vergt worden
geleverd door \cref{ch:b1:integration}, en de weekendopgave hieronder ---
de aangedreven gedempte oscillator --- is het modelgeval voor elke
oscillatievraag in de wetenschappen, van schakelingen tot hangbruggen.
\end{remark}

\section{Oefeningen}

\begin{exercise}[$\star$]\label{exo:b1:diffeq:1}
Los op op $\R$: $\;y' + 2y = \eu^{3x}$; en vervolgens het
\omterm{thm:b1:diffeq:voc}{Cauchy-probleem} $y(0) = 1$.
\end{exercise}

\begin{exercise}[$\star$]\label{exo:b1:diffeq:2}
Los op op $\intoo{0}{+\infty}$: $\;x y' - 2y = x^3$ \emph{(breng de
vergelijking eerst in genormaliseerde vorm)}.
\end{exercise}

\begin{exercise}[$\star$]\label{exo:b1:diffeq:3}
Los op op $\R$ en geef de reële algemene oplossing:
$\;y'' - 3y' + 2y = 0$; $\;y'' + 4y' + 4y = 0$; $\;y'' - 2y' + 5y = 0$.
\end{exercise}

\begin{exercise}[$\star$]\label{exo:b1:diffeq:4}
Los $y'' - y = x^2$ op op $\R$, en vervolgens het \omterm{thm:b1:diffeq:voc}{Cauchy-probleem} $y(0) =
0$, $y'(0) = 1$.
\end{exercise}

\begin{exercise}[$\star\star$]\label{exo:b1:diffeq:5}
Los op op $\intoo{-\frac\pi2}{\frac\pi2}$: $\;y' + y\tan x = \sin 2x$.
\end{exercise}

\begin{exercise}[$\star\star$]\label{exo:b1:diffeq:6}
Los $y'' - 4y' + 3y = (2x + 1)\,\eu^{x}$ op op $\R$. \emph{(Let op de
multipliciteit: is $1$ een wortel van de \omterm{def:b1:diffeq:linear2}{karakteristieke veelterm}?)}
\end{exercise}

\begin{exercise}[$\star\star$]\label{exo:b1:diffeq:7}
Los $y'' + 4y = \sin 2x + x$ op op $\R$ \emph{(superpositie; behandel
elk rechterlid apart)}.
\end{exercise}

\begin{exercise}[$\star\star$]\label{exo:b1:diffeq:8}
Een kop koffie met temperatuur $T_0 = 80\,^\circ$C staat in een kamer van
$20\,^\circ$C. De afkoelingswet van Newton luidt $T' = -k\,(T - 20)$ met
$k > 0$. Los op naar $T(t)$ en bepaal, gegeven dat de koffie na $10$
minuten $50\,^\circ$C is, wanneer ze $25\,^\circ$C bereikt.
\end{exercise}

\begin{exercise}[$\star\star\star$]\label{exo:b1:diffeq:9}
(Gedempte oscillator) Beschouw voor $\varepsilon \geq 0$ de vergelijking
$y'' + 2\varepsilon y' + y = 0$.
\begin{enumerate}
  \item Los op voor $\varepsilon \in \intco{0}{1}$, voor $\varepsilon =
        1$ en voor $\varepsilon > 1$.
  \item Toon aan dat voor $\varepsilon > 0$ elke oplossing naar $0$ gaat
        in $+\infty$, en dat de oplossingen ongelijk aan nul dat voor
        $\varepsilon = 0$ niet doen.
  \item Toon voor $\varepsilon \in \intoo{0}{1}$ aan dat de nulpunten
        van een oplossing ongelijk aan nul regelmatig verdeeld liggen,
        met onderlinge afstand $\frac{\pi}{\sqrt{1 - \varepsilon^2}}$.
\end{enumerate}
\end{exercise}

\begin{exercise}[$\star\star\star$]\label{exo:b1:diffeq:10}
Bepaal alle tweemaal afleidbare functies $f \colon \R \to \R$ met
\[
\forall x, y \in \R, \qquad f(x + y) + f(x - y) = 2 f(x) f(y),
\]
waarbij $f(0) \ne 0$ en $f$ niet constant is.
\emph{Aanwijzing: leg $y$ vast, differentieer tweemaal naar $x$ in $0$;
toon aan dat $f(0) = 1$ en $f'' = c f$ voor een zekere constante $c$;
los daarna op naar het teken van $c$ en ga na welke oplossingen aan de
functievergelijking voldoen.}
\end{exercise}

\begin{exercise}[$\star\star$]\label{exo:b1:diffeq:11}
(Vergelijking van Euler) Los $x^2 y'' - x y' + y = 0$ op op
$\intoo{0}{+\infty}$. \emph{Aanwijzing: zet $z(t) = y(\eu^t)$,
substitueer dus $x = \eu^t$, en toon aan dat $z$ aan een lineaire
vergelijking met constante coëfficiënten voldoet.}
\end{exercise}

\begin{exercise}[$\star\star\star$]\label{exo:b1:diffeq:12}
Beschouw de vergelijking $x\,y' = 2y$ op de hele reële rechte, in de
onbekende afleidbare functie $y \colon \R \to \R$.
\begin{enumerate}
  \item Los op op $\intoo{0}{+\infty}$ en op $\intoo{-\infty}{0}$.
  \item Toon aan dat voor \emph{willekeurige} constanten $a, b \in \R$
        de functie die gelijk is aan $ax^2$ voor $x \geq 0$ en aan
        $bx^2$ voor $x < 0$ afleidbaar is op $\R$ en de vergelijking
        overal oplost.
  \item Besluit dat de oplossingsverzameling op $\R$ een familie met
        twee parameters is, en leg uit waarom dat de uniciteit in
        \cref{thm:b1:diffeq:voc} niet tegenspreekt.
\end{enumerate}
\end{exercise}

\section{Opgave: de aangedreven gedempte oscillator}

\begin{problem}\label{pb:b1:diffeq:1}
Eén vergelijking beheerst een massa aan een veer in een stroperig
medium, de lading in een RLC-kring en een gebouw dat in de wind zwaait:
\[
(E_\Omega)\colon\quad
x'' + 2\lambda x' + \omega_0^2\,x = A\cos(\Omega t),
\]
met $\lambda \geq 0$ de demping, $\omega_0 > 0$ de eigenfrequentie, en
$A > 0$ en $\Omega > 0$ de amplitude en de frequentie van de
aandrijving. Deze opgave haalt het volledige gedrag naar boven: het
uitdoven van het inschakelverschijnsel, het unieke periodieke regime, de
resonantiekromme en haar scherpte (de
\emph{kwaliteitsfactor}), de \omterm{pb:b1:diffeq:1}{zwevingen} in het ongedempte geval, en de
energiebalans die de oscillatie in stand
houdt.\index{resonantie}\index{oscillator}
Tenzij anders vermeld is $0 < \lambda < \omega_0$ (onderkritisch gedempt
regime) en schrijven we $\omega_d = \sqrt{\omega_0^2 - \lambda^2}$.

\medskip
\noindent\textbf{Deel I --- De vrije oscillator.} Hier is $A = 0$.
\begin{enumerate}
  \item Los de homogene vergelijking $(H)$ op voor $0 < \lambda <
        \omega_0$ en voor $\lambda = 0$. (De regimes $\lambda \geq
        \omega_0$ zijn in \cref{exo:b1:diffeq:9} behandeld; citeer ze.)
  \item Toon aan dat voor elke $\lambda > 0$ alle oplossingen van $(H)$
        naar $0$ gaan in $+\infty$ --- in alle drie de regimes.
  \item Definieer de energie $\mathcal E(t) = \frac12 x'(t)^2 +
        \frac12\omega_0^2\,x(t)^2$ langs een oplossing van $(H)$. Toon
        aan dat $\mathcal E'(t) = -2\lambda\,x'(t)^2 \leq 0$, en leid
        daaruit af (zonder iets op te lossen) dat het \omterm{thm:b1:diffeq:voc}{Cauchy-probleem}
        ``$(H)$, $x(t_0) = x'(t_0) = 0$'' alleen de nuloplossing heeft,
        voor elke $\lambda \geq 0$.
  \item Schrijf voor $0 < \lambda < \omega_0$ de oplossing ongelijk aan
        nul als $x(t) = R\,\eu^{-\lambda t}\cos(\omega_d t - \varphi)$
        en zij $T_d = \frac{2\pi}{\omega_d}$ de pseudoperiode. Toon aan
        dat $x(t + T_d) = \eu^{-\lambda T_d}\,x(t)$: elke slingering is
        de vorige, gekrompen met de constante factor $\eu^{-\delta}$
        met $\delta = \frac{2\pi\lambda}{\omega_d}$ (het
        \emph{logaritmische decrement}). Bereken $\delta$ voor
        $\omega_0 = 1$ en $\lambda = 0.1$.
  \item Definieer de \emph{kwaliteitsfactor} $Q =
        \dfrac{\omega_0}{2\lambda}$\index{kwaliteitsfactor}. Toon aan
        dat de oscillator na de tijd $\frac1\lambda$ (één keer de
        amplitude met een factor $\eu$ zien slinken)
        $\frac{\omega_d}{2\pi\lambda}$ pseudoperiodes heeft afgelegd,
        wat bij zwakke demping ($\lambda \ll \omega_0$) ongeveer
        $\frac Q\pi$ is: de kwaliteitsfactor telt, op een factor $\pi$
        na, het aantal oscillaties dat overleeft vóór de amplitude met
        een factor $\eu$ afneemt.
\end{enumerate}

\medskip
\noindent\textbf{Deel II --- Het regime.} Nu is $A > 0$ en $\lambda > 0$.
\begin{enumerate}[resume]
  \item Zoek een particuliere oplossing als reëel deel van
        $z\,\eu^{\iu\Omega t}$ met $z \in \C$
        (\cref{met:b1:diffeq:particular}). Toon aan dat dit werkt met
        \[
        z = \frac{A}{\omega_0^2 - \Omega^2 + 2\iu\lambda\Omega} .
        \]
  \item Leid het regime af in amplitude-fasevorm: $x_p(t) =
        R(\Omega)\cos\bigl(\Omega t - \varphi(\Omega)\bigr)$ met
        \[
        R(\Omega) = \frac{A}{\sqrt{(\omega_0^2 - \Omega^2)^2 +
        4\lambda^2\Omega^2}},
        \qquad
        \tan\varphi = \frac{2\lambda\Omega}{\omega_0^2 -
        \Omega^2},
        \quad \varphi \in \intoo0\pi .
        \]
  \item Duid de twee uiterste regimes: bereken de limieten van $R$ en
        $\varphi$ voor $\Omega \to 0^+$ (quasistatisch antwoord
        $A/\omega_0^2$, fase $0$) en voor $\Omega \to +\infty$ ($R \sim
        A/\Omega^2 \to 0$, fase $\to \pi$: de massa beweegt tegengesteld
        aan een te snelle aandrijving).
  \item Toon aan dat \emph{elke} oplossing van $(E_\Omega)$ gelijk is
        aan $x_p$ plus een oplossing van $(H)$, en dus naar het regime
        $x_p$ convergeert als $t \to +\infty$, wat de beginvoorwaarden
        ook zijn: zodra het inschakelverschijnsel is uitgestorven,
        herinnert de oscillator zich niets meer van zijn start.
  \item Toon aan dat $x_p$ de \emph{enige} periodieke oplossing van
        $(E_\Omega)$ is.
  \item Werk één \omterm{thm:b1:diffeq:voc}{Cauchy-probleem} tot het einde uit: toon voor $x'' + 2x'
        + 2x = \cos t$ met $x(0) = x'(0) = 0$ aan dat de oplossing
        \[
        x(t) = \frac{\cos t + 2\sin t}5
        - \eu^{-t}\,\frac{\cos t + 3\sin t}5
        \]
        is, en wijs het inschakel- en het regimedeel aan.
\end{enumerate}

\medskip
\noindent\textbf{Deel III --- De resonantiekromme.} Studie van $\Omega
\mapsto R(\Omega)$ op $\intoo0{+\infty}$.
\begin{enumerate}[resume]
  \item Zet $u = \Omega^2$ en $g(u) = (\omega_0^2 - u)^2 + 4\lambda^2 u$,
        en toon aan: is $2\lambda^2 < \omega_0^2$, dan bereikt $R$ een
        strikt maximum in de \emph{resonantiefrequentie} $\Omega_r =
        \sqrt{\omega_0^2 - 2\lambda^2}$, met
        \[
        R_{\max} = R(\Omega_r)
        = \frac{A}{2\lambda\sqrt{\omega_0^2 - \lambda^2}} .
        \]
  \item Toon aan dat $\dfrac{R_{\max}}{R(0)} = Q\,\bigl(1 -
        \tfrac{\lambda^2}{\omega_0^2}\bigr)^{-1/2} \geq Q$: in
        \omterm{ex:b1:diffeq:oscillation}{resonantie} wordt de aandrijving (in wezen) met de
        kwaliteitsfactor versterkt.
  \item Bewijs dat de amplitude van de \emph{snelheid} $V(\Omega) =
        \Omega\,R(\Omega)$ precies in $\Omega = \omega_0$ maximaal is
        (en niet in $\Omega_r$), en dat de fase daar $\varphi(\omega_0)
        = \frac\pi2$ is: in $\Omega = \omega_0$ loopt de snelheid
        precies in fase met de kracht.
  \item (Bandbreedte) Los $g(u) = 2\,g(u_r)$ exact op, met $u_r =
        \omega_0^2 - 2\lambda^2$, en leid af dat de twee frequenties
        $\Omega_\pm$ waar $R = R_{\max}/\sqrt2$ voldoen aan $\Omega_+^2
        - \Omega_-^2 = 4\lambda\sqrt{\omega_0^2 - \lambda^2}$; besluit
        dat de bandbreedte bij zwakke demping $\Omega_+ - \Omega_-
        \approx 2\lambda$ is, oftewel $Q \approx
        \frac{\omega_0}{\Omega_+ - \Omega_-}$: scherpe resonantiepieken
        horen bij systemen met hoge $Q$.
  \item Numeriek portret voor $\omega_0 = 1$, $\lambda = 0.05$ ($Q =
        10$) en $A = 1$: bereken $\Omega_r$, $R_{\max}$, het statische
        antwoord $R(0)$ en de benaderde bandbreedte.
  \item Toon aan dat $R$ strikt dalend is op $\intoo0{+\infty}$ zodra
        $2\lambda^2 \geq \omega_0^2$: sterk gedempte systemen hebben in
        het geheel geen resonantiepiek.
\end{enumerate}

\medskip
\noindent\textbf{Deel IV --- Zonder demping: \omterm{pb:b1:diffeq:1}{zwevingen} en \omterm{ex:b1:diffeq:oscillation}{resonantie}.}
Hier is $\lambda = 0$.
\begin{enumerate}[resume]
  \item Bepaal voor $\Omega \neq \omega_0$ de algemene oplossing van
        $x'' + \omega_0^2 x = A\cos(\Omega t)$.
  \item Los het \omterm{thm:b1:diffeq:voc}{Cauchy-probleem} $x(0) = x'(0) = 0$ op en breng het
        antwoord in de productvorm
        \[
        x(t) = \frac{2A}{\omega_0^2 - \Omega^2}\,
        \sin\Bigl(\frac{(\omega_0 - \Omega)t}2\Bigr)
        \sin\Bigl(\frac{(\omega_0 + \Omega)t}2\Bigr) .
        \]
  \item Lees het product voor $\Omega$ dicht bij $\omega_0$ als een
        snelle oscillatie met frequentie $\frac{\omega_0 + \Omega}2$,
        gemoduleerd door een trage omhullende met frequentie
        $\frac{\abs{\omega_0 - \Omega}}2$: de
        \emph{zwevingen}\index{zwevingen}. Geef de periode van de
        omhullende en de maximale amplitude, en merk op hoe beide
        exploderen als $\Omega \to \omega_0$.
  \item Leg $t$ vast en laat $\Omega \to \omega_0$ gaan in de formule
        van vraag 19: toon aan dat de limiet
        \[
        x_\infty(t) = \frac{A\,t\,\sin(\omega_0 t)}{2\omega_0}
        \]
        is, en ga rechtstreeks na dat $x_\infty$ de resonante
        vergelijking $x'' + \omega_0^2 x = A\cos(\omega_0 t)$ oplost met
        $x(0) = x'(0) = 0$ (vergelijk
        \cref{ex:b1:diffeq:oscillation}): \omterm{ex:b1:diffeq:oscillation}{resonantie} is de limiet van
        steeds tragere, steeds grotere \omterm{pb:b1:diffeq:1}{zwevingen}.
  \item Zet de twee lotgevallen van \omterm{ex:b1:diffeq:oscillation}{resonantie} tegenover elkaar:
        lineaire groei $\frac{At}{2\omega_0}$ zonder demping,
        tegenover verzadiging bij $R_{\max} \approx Q\,\frac{A}{\omega_0^2}$
        met zwakke demping. In één zin: welk fysisch mechanisme zet het
        eerste om in het tweede?
\end{enumerate}

\medskip
\noindent\textbf{Deel V --- Energiebalans en synthese.}
\begin{enumerate}[resume]
  \item Bereken in het regime van Deel II het gemiddelde over één
        periode $\frac{2\pi}\Omega$ van (a) het vermogen dat de
        aandrijving levert, $P_{\mathrm{in}}(t) = A\cos(\Omega t) \cdot
        x_p'(t)$, en (b) het vermogen dat de demping verbruikt,
        $P_{\mathrm{diss}}(t) = 2\lambda\,x_p'(t)^2$. Toon aan dat
        beide gemiddelden gelijk zijn aan $\lambda\,R^2\Omega^2$: de
        aandrijving voert precies aan wat de demping opstookt --- en
        daarom is het regime stationair.
  \item Waar precies gebruikte de opgave: (i) de structuurstelling
        \cref{thm:b1:diffeq:structure2}; (ii) de methode van de
        complexe exponentiële; (iii) een functiestudie in reële
        variabele in de stijl van \cref{ch:b1:functions}? Eén zin per
        onderdeel.
  \item Synthese: beschrijf het volledige gedrag van $(E_\Omega)$ ---
        vrij tegenover aangedreven, gedempt tegenover ongedempt, de rol
        van $Q$ als de ene dimensieloze knop die piekhoogte,
        bandbreedte en levensduur van het inschakelverschijnsel
        afstemt --- en vermeld waar het verhaal verdergaat: $2 \times
        2$-stelsels van de eerste orde (\cref{ch:b1:matrices} en het
        volume van bachelorjaar 2) en de ontbinding van een algemene
        periodieke aandrijving in sinusoïden (fourierreeksen, in het
        volume van bachelorjaar 3), waarvoor het sinusoïdale geval van
        deze opgave de fundamentele bouwsteen is.
\end{enumerate}
\end{problem}
